% Options for packages loaded elsewhere
\PassOptionsToPackage{unicode}{hyperref}
\PassOptionsToPackage{hyphens}{url}
\documentclass[
]{article}
\usepackage{xcolor}
\usepackage[height=11in,width=8.5in,top=3cm,bottom=3cm,left=2cm,right=2cm,headsep=10pt,letterpaper]{geometry}
\usepackage{amsmath,amssymb}
\setcounter{secnumdepth}{5}
\usepackage{iftex}
\ifPDFTeX
  \usepackage[T1]{fontenc}
  \usepackage[utf8]{inputenc}
  \usepackage{textcomp} % provide euro and other symbols
\else % if luatex or xetex
  \usepackage{unicode-math} % this also loads fontspec
  \defaultfontfeatures{Scale=MatchLowercase}
  \defaultfontfeatures[\rmfamily]{Ligatures=TeX,Scale=1}
\fi
\usepackage{lmodern}
\ifPDFTeX\else
  % xetex/luatex font selection
\fi
% Use upquote if available, for straight quotes in verbatim environments
\IfFileExists{upquote.sty}{\usepackage{upquote}}{}
\IfFileExists{microtype.sty}{% use microtype if available
  \usepackage[]{microtype}
  \UseMicrotypeSet[protrusion]{basicmath} % disable protrusion for tt fonts
}{}
\makeatletter
\@ifundefined{KOMAClassName}{% if non-KOMA class
  \IfFileExists{parskip.sty}{%
    \usepackage{parskip}
  }{% else
    \setlength{\parindent}{0pt}
    \setlength{\parskip}{6pt plus 2pt minus 1pt}}
}{% if KOMA class
  \KOMAoptions{parskip=half}}
\makeatother
\usepackage{longtable,booktabs,array}
\usepackage{calc} % for calculating minipage widths
% Correct order of tables after \paragraph or \subparagraph
\usepackage{etoolbox}
\makeatletter
\patchcmd\longtable{\par}{\if@noskipsec\mbox{}\fi\par}{}{}
\makeatother
% Allow footnotes in longtable head/foot
\IfFileExists{footnotehyper.sty}{\usepackage{footnotehyper}}{\usepackage{footnote}}
\makesavenoteenv{longtable}
\usepackage{graphicx}
\makeatletter
\newsavebox\pandoc@box
\newcommand*\pandocbounded[1]{% scales image to fit in text height/width
  \sbox\pandoc@box{#1}%
  \Gscale@div\@tempa{\textheight}{\dimexpr\ht\pandoc@box+\dp\pandoc@box\relax}%
  \Gscale@div\@tempb{\linewidth}{\wd\pandoc@box}%
  \ifdim\@tempb\p@<\@tempa\p@\let\@tempa\@tempb\fi% select the smaller of both
  \ifdim\@tempa\p@<\p@\scalebox{\@tempa}{\usebox\pandoc@box}%
  \else\usebox{\pandoc@box}%
  \fi%
}
% Set default figure placement to htbp
\def\fps@figure{htbp}
\makeatother
% definitions for citeproc citations
\NewDocumentCommand\citeproctext{}{}
\NewDocumentCommand\citeproc{mm}{%
  \begingroup\def\citeproctext{#2}\cite{#1}\endgroup}
\makeatletter
 % allow citations to break across lines
 \let\@cite@ofmt\@firstofone
 % avoid brackets around text for \cite:
 \def\@biblabel#1{}
 \def\@cite#1#2{{#1\if@tempswa , #2\fi}}
\makeatother
\newlength{\cslhangindent}
\setlength{\cslhangindent}{1.5em}
\newlength{\csllabelwidth}
\setlength{\csllabelwidth}{3em}
\newenvironment{CSLReferences}[2] % #1 hanging-indent, #2 entry-spacing
 {\begin{list}{}{%
  \setlength{\itemindent}{0pt}
  \setlength{\leftmargin}{0pt}
  \setlength{\parsep}{0pt}
  % turn on hanging indent if param 1 is 1
  \ifodd #1
   \setlength{\leftmargin}{\cslhangindent}
   \setlength{\itemindent}{-1\cslhangindent}
  \fi
  % set entry spacing
  \setlength{\itemsep}{#2\baselineskip}}}
 {\end{list}}
\usepackage{calc}
\newcommand{\CSLBlock}[1]{\hfill\break\parbox[t]{\linewidth}{\strut\ignorespaces#1\strut}}
\newcommand{\CSLLeftMargin}[1]{\parbox[t]{\csllabelwidth}{\strut#1\strut}}
\newcommand{\CSLRightInline}[1]{\parbox[t]{\linewidth - \csllabelwidth}{\strut#1\strut}}
\newcommand{\CSLIndent}[1]{\hspace{\cslhangindent}#1}
\setlength{\emergencystretch}{3em} % prevent overfull lines
\providecommand{\tightlist}{%
  \setlength{\itemsep}{0pt}\setlength{\parskip}{0pt}}
\makeatletter
\@ifpackageloaded{subfig}{}{\usepackage{subfig}}
\@ifpackageloaded{caption}{}{\usepackage{caption}}
\captionsetup[subfloat]{margin=0.5em}
\AtBeginDocument{%
\renewcommand*\figurename{Figure}
\renewcommand*\tablename{Table}
}
\AtBeginDocument{%
\renewcommand*\listfigurename{List of Figures}
\renewcommand*\listtablename{List of Tables}
}
\newcounter{pandoccrossref@subfigures@footnote@counter}
\newenvironment{pandoccrossrefsubfigures}{%
\setcounter{pandoccrossref@subfigures@footnote@counter}{0}
\begin{figure}\centering%
\gdef\global@pandoccrossref@subfigures@footnotes{}%
\DeclareRobustCommand{\footnote}[1]{\footnotemark%
\stepcounter{pandoccrossref@subfigures@footnote@counter}%
\ifx\global@pandoccrossref@subfigures@footnotes\empty%
\gdef\global@pandoccrossref@subfigures@footnotes{{##1}}%
\else%
\g@addto@macro\global@pandoccrossref@subfigures@footnotes{, {##1}}%
\fi}}%
{\end{figure}%
\addtocounter{footnote}{-\value{pandoccrossref@subfigures@footnote@counter}}
\@for\f:=\global@pandoccrossref@subfigures@footnotes\do{\stepcounter{footnote}\footnotetext{\f}}%
\gdef\global@pandoccrossref@subfigures@footnotes{}}
\@ifpackageloaded{float}{}{\usepackage{float}}
\floatstyle{ruled}
\@ifundefined{c@chapter}{\newfloat{codelisting}{h}{lop}}{\newfloat{codelisting}{h}{lop}[chapter]}
\floatname{codelisting}{Listing}
\newcommand*\listoflistings{\listof{codelisting}{List of Listings}}
\makeatother
\usepackage{bookmark}
\IfFileExists{xurl.sty}{\usepackage{xurl}}{} % add URL line breaks if available
\urlstyle{same}
\hypersetup{
  pdftitle={Synaptic Plasticity Model Predicts Optimum Treatment Parameters for Amblyopia},
  hidelinks,
  pdfcreator={LaTeX via pandoc}}

\title{Synaptic Plasticity Model Predicts Optimum Treatment Parameters
for Amblyopia}
\author{}
\date{}

\begin{document}
\maketitle

\section{Main Manuscript for}\label{sec:main-manuscript-for}

\textbf{Synaptic Plasticity Model Predicts Optimum Treatment Parameters
for Amblyopia}

Brian S. Blais\(^{1*}\) and Eric Gaier\(^{2,3}\)

\(^{1}\)Department of Biological and Biomedical Sciences, Bryant
University

\(^{2}\)Picower Institute for Learning and Memory, Massachusetts
Institute of Technology

\(^{3}\)Department of Ophthalmology, Boston Children's Hospital, Harvard
Medical School

\(^{*}\)Brian Blais corresponding author

\textbf{Email}: bblais@bryant.edu

\textbf{Author Contributions}: Paste the author contributions here.

\textbf{Competing Interest Statement}: No competing interests.

\textbf{Classification}: Biological Sciences - Neuroscience

\textbf{Keywords}: amblyopia, synaptic plasticity, vision

\textbf{This PDF file includes}: - Main Text - Figures 1 to X - Tables 1
to X

\subsection{Abstract}\label{sec:abstract}

\begin{itemize}
\tightlist
\item[$\square$]
  Purpose: Extend a biologically motivated synaptic plasticity model
  (BCM with natural image input) to simulate common amblyopia
  treatments.
\item[$\square$]
  Methods: Use feedforward-only architecture in PlasticNet to implement
  patching, atropine penalization, dichoptic training, and contrast
  balancing protocols.
\item[$\square$]
  Findings: Predicts treatment-specific recovery trajectories and
  identifies optimal stimulus parameters (e.g., contrast levels, patch
  durations).
\item[$\square$]
  Implications: Offers quantitative insight into the timing and dosing
  of interventions, with predictions aligning with clinical results.
\end{itemize}

Amblyopia is a common cause of visual impairment that results from
unequal visual inputs during development, known to manifest through
synaptic alterations in the visual cortex. What is not known is the
detailed mechanisms of these synaptic changes and how these mechanisms
impact the dynamics of recovery. Here we use a computational model of
neural plasticity, the Bienenstock, Cooper, and Munro (BCM) model
(Bienenstock et al 1982), to compare the dynamics of amblyopia recovery
at the neuronal level under several treatment protocols, including
optical correction, patching, atropine penalization, and binocular
therapies. We use this model to determine optimal parameters for the
treatment of amblyopia, showing that recovery achieved with dichoptic
masks combined with an interocular contrast disparity exceeded that of
patch and atropine treatments. Further, patch and atropine treatment
models produced faster recovery compared to a contrast disparity alone,
highlighting the importance of the dichoptic masks. The rate of recovery
depended on treatment features such as the size of the dichoptic masks
and the magnitude of the contrast disparity, both experimentally
accessible. In this way, the model suggests optimal values for these
modifications.

\textbf{Significance Statement}

\begin{quote}
Paste your significance statement here. Please note that it should not
exceed 120 words, but should be at least 50 words in length. It should
not include any references.
\end{quote}

\section{Main Text}\label{sec:main-text}

\subsection{Introduction}\label{sec:introduction}

\begin{itemize}
\item[$\square$]
  Brief review of amblyopia development and your prior findings (cite
  previous paper).
\item[$\square$]
  Rationale: Most treatments aim to rebalance interocular competition,
  but empirical optimization is challenging.
\item[$\square$]
  Common treatments:

  \begin{itemize}
  \tightlist
  \item[$\square$]
    Patching: Occlusion of dominant eye.
  \item[$\square$]
    Atropine: Blurring the dominant eye pharmacologically.
  \item[$\square$]
    Dichoptic training: Presenting stimuli separately to each eye to
    promote binocular integration.
  \item[$\square$]
    Contrast balancing (e.g., Luminopia): Reduce contrast to dominant
    eye in a movie/game format.
  \end{itemize}
\item[$\square$]
  Goal: Use a synaptic plasticity model to simulate these treatments and
  explore parameter sensitivities.
\item[$\square$]
  Hypothesis: Optimal treatments balance synaptic competition without
  destabilizing cortical representations.
\item[$\square$]
  BCM learning and recovery:

  \begin{itemize}
  \tightlist
  \item[$\square$]
    Homeostatic and competitive nature is ideal for simulating
    rebalancing.
  \item[$\square$]
    Prior modeling of amblyopia treatment: Summary of efforts in
    literature.
  \end{itemize}
\item[$\square$]
  PlasticNet advantages:

  \begin{itemize}
  \tightlist
  \item[$\square$]
    Realistic inputs (natural images)
  \item[$\square$]
    Parameter flexibility for fine-grained treatment modeling
  \end{itemize}
\item[$\square$]
  Clinical motivation: Need for tailored treatments by severity, age,
  and amblyopia type. \#\# Results
\item[$\square$]
  Baseline (no treatment): Persistent suppression of amblyopic eye.
\item[$\square$]
  Patching:

  \begin{itemize}
  \tightlist
  \item[$\square$]
    Effective in early periods; risk of reverse amblyopia with excessive
    patching.
  \end{itemize}
\item[$\square$]
  Atropine:

  \begin{itemize}
  \tightlist
  \item[$\square$]
    Mimics patching effect with softer recovery dynamics.
  \end{itemize}
\item[$\square$]
  Dichoptic training:

  \begin{itemize}
  \tightlist
  \item[$\square$]
    Strongest recovery in binocular metrics when balanced properly.
  \item[$\square$]
    Sensitive to contrast parameters.
  \item[$\square$]
    Contrast reduction (Luminopia-like):
  \item[$\square$]
    Continuous contrast adaptation of dominant eye yields gradual OD
    recovery.
  \end{itemize}
\item[$\square$]
  Optimal parameters:

  \begin{itemize}
  \tightlist
  \item[$\square$]
    Identified sweet spots for contrast ratio and patching time per
    treatment type.
  \item[$\square$]
    Early intervention window (plasticity-dependent) strongly shapes
    outcome.
  \end{itemize}
\item[$\square$]
  Include:

  \begin{itemize}
  \tightlist
  \item[$\square$]
    OD histogram evolution over training.
  \item[$\square$]
    Time-course plots of synaptic weight recovery.
  \item[$\square$]
    Heatmaps of treatment efficacy by parameter.
  \end{itemize}
\end{itemize}

\subsection{Discussion}\label{sec:discussion}

\begin{itemize}
\tightlist
\item[$\square$]
  Key insight: Model confirms and refines clinical understanding of
  treatment dynamics.
\item[$\square$]
  Reaffirms sensitive period effects.
\item[$\square$]
  Shows contrast balancing is more forgiving than patching in timing.
\item[$\square$]
  Comparison with human data:

  \begin{itemize}
  \tightlist
  \item[$\square$]
    Atropine vs.~patching trials (e.g., PEDIG).
  \item[$\square$]
    Recent dichoptic training and Luminopia results.
  \end{itemize}
\item[$\square$]
  Treatment optimization:

  \begin{itemize}
  \tightlist
  \item[$\square$]
    Parameter tuning can guide personalized treatment plans.
  \end{itemize}
\item[$\square$]
  Limitations:

  \begin{itemize}
  \tightlist
  \item[$\square$]
    Feedforward only: no attentional or reward mechanisms modeled.
  \item[$\square$]
    No behavioral feedback loop.
  \end{itemize}
\item[$\square$]
  Future work:

  \begin{itemize}
  \tightlist
  \item[$\square$]
    Combine with reinforcement learning agent to model perceptual
    learning tasks.
  \item[$\square$]
    Add spatial structure (e.g., orientation maps) for richer V1
    modeling.
  \item[$\square$]
    Apply to adult treatment paradigms (model reduced plasticity).
  \end{itemize}
\item[$\square$]
  BCM-based model predicts differing recovery dynamics under common
  amblyopia treatments.
\item[$\square$]
  Early, contrast-based interventions may provide high efficacy with low
  risk.
\item[$\square$]
  Synaptic plasticity modeling offers a promising tool for treatment
  planning and hypothesis generation.
\end{itemize}

\subsection{Materials and Methods}\label{sec:materials-and-methods}

\begin{itemize}
\tightlist
\item[$\square$]
  Model framework:
\item[$\square$]
  Same BCM-based PlasticNet model as prior paper.
\item[$\square$]
  Same strabismus and/or anisometropia manipulations used to induce
  amblyopia.
\item[$\square$]
  Treatment protocols modeled:

  \begin{itemize}
  \tightlist
  \item[$\square$]
    Patching: One eye's input occluded (removed) for a percentage of
    training epochs.
  \item[$\square$]
    Atropine: One eye blurred (e.g., Gaussian filter) simulating
    low-acuity vision.
  \item[$\square$]
    Dichoptic: Present spatially distinct images to each eye with
    different contrast scaling.
  \item[$\square$]
    Contrast balancing: Gradually or abruptly reduce contrast to the
    dominant eye during image input.
  \end{itemize}
\item[$\square$]
  Evaluation metrics:

  \begin{itemize}
  \tightlist
  \item[$\square$]
    Recovery time to balanced ocular dominance.
  \item[$\square$]
    Receptive field restoration (sharpness, binocular selectivity).
  \end{itemize}
\item[$\square$]
  Dependence on treatment onset timing and duration.
\item[$\square$]
  Explored parameters:

  \begin{itemize}
  \tightlist
  \item[$\square$]
    \% patching time
  \item[$\square$]
    Contrast ratios
  \item[$\square$]
    Timing of treatment onset (early vs.~late)
  \end{itemize}
\end{itemize}

\subsection{Acknowledgments}\label{sec:acknowledgments}

Paste your acknowledgments here.

\subsection{References}\label{sec:references}

\begin{itemize}
\tightlist
\item[$\square$]
  Include:

  \begin{itemize}
  \tightlist
  \item[$\square$]
    Clinical trials (PEDIG studies, Luminopia trials).
  \item[$\square$]
    Models of binocular recovery.
  \item[$\square$]
    Synaptic plasticity and BCM literature.
  \end{itemize}
\end{itemize}

\protect\phantomsection\label{refs}
\begin{CSLReferences}{0}{1}
\bibitem[\citeproctext]{ref-glaser2002randomized}
\CSLLeftMargin{1. }%
\CSLRightInline{Glaser, S.R., Matazinski, A.M., Sclar, D.M., Sala, N.A.,
Vroman, C.M., Tanner, C.E., Stager, D.R., Berry, P.M., Felius, J.,
Wilkerson, J.A., et al. (2002). A randomized trial of atropine vs
patching for treatment of moderate amblyopia in children. Archives of
ophthalmology \emph{120}, 268--278.}

\end{CSLReferences}

\subsection{Figures and Tables}\label{sec:figures-and-tables}

\begin{quote}
This section actually seems a bit superfluous right now, I wonder if we
don't need to try and link the ocular dominance measures with visual
acuity

Instead, we could focus the conclusion on linking the directional
conclusions with ODI/day to the existing experimental literature, and
making recommendations for future amblyopia treatment studies. Thoughts?
\end{quote}

Now that we have a system of simulation environments to explore, we can
compare to experimentally observed rates of recovery.
From\textsuperscript{1} we have results from several visual protocols.

\begin{enumerate}
\def\labelenumi{\arabic{enumi}.}
\tightlist
\item
  Only those patients are included if they had their \emph{refractive
  error corrected for at least 4 weeks}
\item
  In the patching group most patients received \emph{no more than 6-8
  hours of patching per day}
\item
  The resulting improvement in the visual acuity (measured in lines) is
  given here:
\end{enumerate}

\begin{longtable}[]{@{}lll@{}}
\toprule\noalign{}
Time & Patch {[}lines{]} & Atropine {[}lines{]} \\
\midrule\noalign{}
\endhead
\bottomrule\noalign{}
\endlastfoot
5 weeks & \(+2.22 \pm 0.2\) & \(+1.37 \pm 0.2\) \\
16 weeks & \(+2.94 \pm 0.2\) & \(+2.42 \pm 0.2\) \\
24 weeks & \(+3.16 \pm 0.2\) & \(+2.84 \pm 0.2\) \\
\end{longtable}

This small amount of data lets us estimate the relative rates of
improvement from the treatments. Since the patch treatment is only about
1/3 day, the total time for treatment would be
\(19 \text{weeks}\times \frac{7 \text{day}}{1 \text{week}}\times 1/3=44 \text{day}\)
For patch treatment with the above data we have a rate of about
\(0.94 \text{lines} / 44 \text{day}=0.021 \text{lines}/\text{day}\).
Likewise, for atropine, we have a rate of about
\(1.47\text{lines} / 133 \text{day}=0.011 \text{lines}/\text{day}\). So
the patch treatment is approximately twice as fast as the atropine.
Looking at Figure Figure~\ref{fig:dODI_atropine_vs_blur} we see that
this can put a rough constraint on the parameters. For a closed-eye
noise for the patch treatment of \(\sigma_n=0.8\) (recovery rate ODI/day
\(\sim 0.2\)), the atropine treatment must have a lower noise level --
we can look at the atropine parameters which yield recover rates ODI/day
\(\sim 0.1\)). For little blur, we need a noise level of around
\(\sigma_n=0.6\), but if the atropine produces a significant blur, then
the noise level of those inputs must be much lower -- well below
\(\sigma_n=0.3\) for blur filter size 6.0, for example.

This noise level for atropine is entirely consistent with the same
open-eye noise level with the glasses ``fix'' discussed earlier. Here we
have an independent line of argument to suggest that atropine may blur
the natural input, but doesn't change the overall spontaneous activity
of neurons. Further, it suggest that there is a significant
physiological different in the activity distributions between
unstructure input (e.g.~patch) and degraded input (e.g.~atropine).

In this way we may hope to constrain other parameters of the model by
comparing to experimental rates of recovery.

\subsubsection{Figures and Tables}\label{sec:figures-and-tables-1}

\begin{figure}
\centering
\pandocbounded{\includegraphics[keepaspectratio,alt={Recovery rate (ocular dominance shift per time) for simulated treatments of amblyopia. The simulations show that for a wide range of contrasts (10\%-50\%) and mask sizes, the binocular treatment is comparable to the patch and glasses treatment. For a narrow range of contrasts (around 30-40\%) and small overlap with the masks, the binocular treatment is superior to the patch and glasses treatments.}]{/Users/bblais/Documents/Git/Amblyopia-Simulation/Manuscript/resources/Pasted image 20250416110214.png}}
\caption{Recovery rate (ocular dominance shift per time) for simulated
treatments of amblyopia. The simulations show that for a wide range of
contrasts (10\%-50\%) and mask sizes, the binocular treatment is
comparable to the patch and glasses treatment. For a narrow range of
contrasts (around 30-40\%) and small overlap with the masks, the
binocular treatment is superior to the patch and glasses
treatments.}\label{fig:dODI_contrast}
\end{figure}

\#todo

\begin{itemize}
\tightlist
\item[$\square$]
  ORCID number attached to the PNAS profile for each author
\item[$\square$]
  Classifications at
  \url{https://www.pnas.org/author-center/submitting-your-manuscript}

  \begin{itemize}
  \tightlist
  \item
    Biological Sciences -- Neuroscience
  \item
    Biological Sciences - Biophysics and Computational Biology
  \end{itemize}
\item
  {[} {]}
\end{itemize}

\end{document}
